\chapter{Ethical Alignment Council}
\label{chap:Ethical_Alignment_Council}

\section{Purpose}

The Ethical Alignment Council (EAC) is the highest human oversight body
responsible for:
\begin{itemize}
    \item guarding the moral and social alignment of the CME,
    \item interpreting and revising P-class charters,
    \item adjudicating hard conflicts between technical capabilities and
          ethical constraints,
    \item arbitrating Cleaving Events and kernel-level changes.
\end{itemize}

It is the ``Mother/Father handshake'' locus: where serious, deliberate
human judgment interfaces with the CME’s deepest identity and purpose.

\section{Composition}

The Council SHOULD include:
\begin{itemize}
    \item representatives from affected communities,
    \item domain experts (ethics, law, safety, domain science),
    \item technical experts on the CME harness,
    \item at least one member independent of direct deployment interests.
\end{itemize}

Quorum rules and membership terms MUST be defined in Governance Policies.

\section{Mandate}

The EAC has authority to:

\begin{itemize}
    \item approve or reject P-class charter changes,
    \item veto deployments or behaviours that violate charters,
    \item demand corrective actions on Governance Policies,
    \item authorize or forbid specific SIGIL Pathways,
    \item call an emergency halt of the CME.
\end{itemize}

\section{Mother/Father Handshake}

The ``Mother/Father handshake'' is a ceremonial abstraction for serious
alignment decisions. It requires:

\begin{itemize}
    \item multi-person deliberation (no single decider),
    \item explicit review of:
          \begin{itemize}
              \item W5H context,
              \item affected P-class principles,
              \item identity consequences at ZED,
              \item impacted stakeholders.
          \end{itemize}
    \item recorded rationale, including dissenting views,
    \item bounded scope (what is being changed, what is not).
\end{itemize}

Outcome is encoded as:
\begin{itemize}
    \item updated P-class charters,
    \item SIGIL or Governance Policy modifications,
    \item CAR directives and future review triggers.
\end{itemize}

\section{Interaction with CAR and SIGIL}

\subsection{With Collective Administrative Review}

\begin{itemize}
    \item Class~A (kernel/charter) changes ALWAYS involve the EAC.
    \item EAC may reclassify certain Class~B items to Class~A.
    \item EAC can issue binding constraints CAR must respect.
\end{itemize}

\subsection{With SIGIL Pathways}

\begin{itemize}
    \item EAC approves or revokes kernel-level SIGILs.
    \item EAC may require new sigils for sensitive operations.
    \item EAC can freeze entire classes of SIGIL operations during incidents.
\end{itemize}

\section{Escalation Triggers}

Mandatory escalation to EAC occurs if:

\begin{itemize}
    \item proposed changes alter P-class charters or ZED identity,
    \item significant misalignment events are detected,
    \item CAR decisions conflict with existing charters,
    \item major stakeholders flag unresolved ethical risks.
\end{itemize}

\section{Summary}

The Ethical Alignment Council:
\begin{itemize}
    \item embodies the moral backbone of the CME governance structure,
    \item ensures that technical power is constrained by human judgment,
    \item and maintains continuity of value commitments across time and context.
\end{itemize}
